\section{Preorder, Inorder, and Postorder in One Traversal}
\label{sec:trees-all-three}

\problemstatement{Compute all three depth-first traversals -- preorder,
inorder, and postorder -- in a \emph{single} pass over the tree, using one
explicit stack.}

\begin{patternbox}
\emph{``All three in one traversal''} is a state-machine problem. Each node
is genuinely visited three times during a depth-first walk: on the way
down (before its left subtree), between its subtrees, and on the way up
(after its right subtree). Attaching a \textbf{state counter $1/2/3$} to
each stacked node lets one loop record the node into the right list at the
right moment.
\end{patternbox}

\subsection*{Understanding the problem carefully}

Picture the classic ``draw a line around the tree'' visualisation of a
DFS: your pen touches each node three times -- once on its left, once
underneath, once on its right. Preorder records the first touch, inorder
the second, postorder the third. If we stack each node together with a
counter saying which touch is next, a single loop can dispatch it to the
correct output list and advance its counter.

\begin{ideabox}[One Node, Three Touches, Three States]
Push each node with state $1$. When popped:
\begin{itemize}
  \item \textbf{State 1} (first touch): add to \emph{preorder}; increment
        its state to $2$ and re-push it; then push its left child (state
        $1$).
  \item \textbf{State 2} (second touch): add to \emph{inorder}; increment
        to $3$ and re-push; then push its right child (state $1$).
  \item \textbf{State 3} (third touch): add to \emph{postorder}; discard
        the node.
\end{itemize}
\end{ideabox}

\subsection*{Algorithm}

\begin{enumerate}
  \item Push $(\text{root}, 1)$ onto a stack. While non-empty, pop
        $(node, state)$:
  \begin{enumerate}
    \item If $state=1$: append to preorder; push $(node,2)$; if left
          child exists push $(\text{left},1)$.
    \item If $state=2$: append to inorder; push $(node,3)$; if right
          child exists push $(\text{right},1)$.
    \item If $state=3$: append to postorder.
  \end{enumerate}
\end{enumerate}

\subsection*{Dry run}

\dryrun{Tree: root $1$; left $2$; right $3$.}

\begin{itemize}
  \item Pop $(1,1)$: pre$=[1]$; push $(1,2)$; push $(2,1)$.
  \item Pop $(2,1)$: pre$=[1,2]$; push $(2,2)$ (no left).
  \item Pop $(2,2)$: in$=[2]$; push $(2,3)$ (no right).
  \item Pop $(2,3)$: post$=[2]$.
  \item Pop $(1,2)$: in$=[2,1]$; push $(1,3)$; push $(3,1)$.
  \item Pop $(3,1)$: pre$=[1,2,3]$; push $(3,2)$. Pop $(3,2)$:
        in$=[2,1,3]$; push $(3,3)$. Pop $(3,3)$: post$=[2,3]$.
  \item Pop $(1,3)$: post$=[2,3,1]$.
  \item Pre $=1,2,3$; \; In $=2,1,3$; \; Post $=2,3,1$.
\end{itemize}

\subsection*{C++ implementation}

\begin{lstlisting}
#include <bits/stdc++.h>
using namespace std;

struct TreeNode {
    int val;
    TreeNode* left;
    TreeNode* right;
    TreeNode(int v) : val(v), left(nullptr), right(nullptr) {}
};

void allThree(TreeNode* root, vector<int>& pre,
              vector<int>& in, vector<int>& post) {
    if (root == nullptr) return;
    stack<pair<TreeNode*, int>> st;
    st.push({root, 1});

    while (!st.empty()) {
        auto [node, state] = st.top();
        st.pop();
        if (state == 1) {
            pre.push_back(node->val);
            st.push({node, 2});
            if (node->left) st.push({node->left, 1});
        } else if (state == 2) {
            in.push_back(node->val);
            st.push({node, 3});
            if (node->right) st.push({node->right, 1});
        } else {
            post.push_back(node->val);
        }
    }
}

int main() {
    TreeNode* root = new TreeNode(1);
    root->left = new TreeNode(2);
    root->right = new TreeNode(3);

    vector<int> pre, in, post;
    allThree(root, pre, in, post);

    auto printv = [](const string& n, vector<int>& v) {
        cout << n << ": ";
        for (int x : v) cout << x << " ";
        cout << "\n";
    };
    printv("pre ", pre);
    printv("in  ", in);
    printv("post", post);
    return 0;
}
\end{lstlisting}

\begin{complexitybox}
\textbf{Time Complexity.} $O(n)$ -- each node is pushed and popped exactly
three times, a constant factor. \textbf{Space Complexity.} $O(h)$ for the
stack, plus $O(n)$ for the three output arrays.
\end{complexitybox}

\begin{takeawaybox}
The state $1/2/3$ per stack frame is exactly what a recursive DFS stores
implicitly as ``which line am I on when I return.'' Making that state
explicit is a general technique for merging several traversal-order tasks
into one pass.
\end{takeawaybox}
